% =============================================================================
% CHAPTER 2: Background
% =============================================================================
% SPDX-FileCopyrightText: 2026 Frank Langbein <frank@langbein.org>
% SPDX-License-Identifier: LPPL-1.3c
%
% Suggested sections: Context (or Wider context / Domain and context); Related
% work (or Literature review / Prior work); Problem statement (and research
% question(s)). Optional: Theory and concepts; Stakeholders and constraints
% (as subsections of Context or separate). End with a clear problem statement.
% Assume a competent reader; avoid lengthy descriptions of standard tools.
% =============================================================================

\section{Context}\label{sec:context}

Cyber security education has steadily moved towards practical scenario based learning, honing skills rather than declaring them through guided courses. Students need to be able to understand how specific vulnerabilities are exploited and work 'in the wild' rather than just understanding definitions such as SQL injection or privilege escalation. \gls{ctf} exercises have become a widely used educational tool, with prior work arguing that they complement traditional teaching and cover knowledge and skills in a way that map well to cyber security training (Švábenský et al., 2021)\cite{svabensky_automated_2024}. Within this, VulnForge attempts to address a specific aspect of this problem, whether generative AI can be used to create and manage personalised vulnerable training environments that remain educationally useful, address shortcomings of existing applications, operationally safe and practically deployable.

Existing cyber security training is dominated by static, pre-authored challenges that will often expose the learner to the same vulnerability, scenario, and flag each time. There is no mainstream platform that currently offers generative AI to synthesis vulnerabilities paired with a context aware tutor. Providing such a platform can help turn abstract concepts into simulated practice. Commercial and community platforms were built with a shared architectural assumption; once a challenge is authored, reviewed, and published, every user will encounter an identical artifact. Although this assumption provides positive operational economics, it also underlies the most commonly cited pedagogical weakness. The static challenges tend to overemphasize technical exploitation, without the user being able to explore limitation of their understanding. Take a web vulnerability, some of the web platforms may tell the user to craft a command to exploit a path traversal\cite{noauthor_cwe_nodate}, however, the user hasn't been able to yet use their knowledge of security to realise that is what they needed. 

The demand for these cyber security skills has grown considerably faster than education around them. The 2024 ISC2 Cybersecurity Workforce study estimated that the global workforce gap in relation to professionals needed to defend organisations reached 4.7 million in 2024, a 19.1 percent increase from 2023\cite{noauthor_results_nodate}. The study also reported that 30 percent of respondents reported skills gaps within their teams, and 33 percent reported staffing shortages. In addition, the United States Bureau of Labor Statistics projects a 29 percent growth rate for information security analysts through to 2034\cite{noauthor_information_nodate}, demonstrating that this is not merely a hiring problem, there is a demonstrable skill gap. 

\section{Related work}\label{sec:related}

Existing platforms have already proven the value of hands on training. Hack The Box Academy has a traditional CTF route. It presents cybersecurity training in a highly interactive process, with target hosts, checkpoints, exercises and skill assessments that allow learners to reproduce techniques they have learnt about rather than simply read about them. Katsantonis et al. 2023 \cite{katsantonis_cyber_2023} argues that cyber ranges have become important learning and training tools, but notes that their design is often hindered by a lack of shared standards and methodologies. VulnForge is not just attempting to generate vulnerable code, but rather it has a broader system, having an LLM produce structured vulnerability specifications, an orchestration layer to deploy these in isolated challenge environments, allowing the learner to access them remotely and providing a tutor to give contextual hints. The proposed solution is closer to a cyber range than a regular CTF. Lazarov et al. (2025)\cite{lazarov_lessons_2025} reinforces this idea by describing cyber ranges as controlled environments for deploying virtual machines and networks, and by having increased interactivity (the contextual tutor in this case), it can improve engagement and the development of practical skills. Therefore, the educational requirement is not only creating exploitable targets, but also environments in a system that supports learning. 

TryHackMe\cite{tryhackme} (fix reference) founded in 2018, takes a beginner oriented approach, organising content into rooms, combining guided labs with written instructions with deployable targets. Similar to Hack The Box, each deployed instance contains identical vulnerabilities and accepts the same flag, with exceptions for King of the Hill modes that rotate flags hourly whilst leaving the vulnerability unchanged. picoCTF\cite{picoctf}, the largest free student-oriented CTF also follows a similar model, cloning vulnerable web applications for each learner, with a lack of content variation. 

Including aspects of a CTF is particularly relevant here as it provides a manageable pedagogical unit, bounded with a clear task and goal, with observable progress and explicit proof of successful exploitation. Švábenský et al. (2021)\cite{svabensky_automated_2024} describes CTFs as a popular form of modern cyber security education, having flags as a convenient proof of exploitation, with challenge categories that are mapped to a specific class demonstrates the appeal of this model. Traditional CTF content does raise questions about replay-ability, variety and personalisation, so once the users have exhausted a fixed content library, the educational value rapidly diminishes. 

\section{Procedural Challenge Generation}\label{sec:procedural}
There has been some work where existing systems have attemtped to anticipate this weakness, and have some variety to CTF content. The fou

\section{LLMs as authors of insecure code}\label{sec:LLMs}
OpenAI's Codex release in 2023 established that transformer LLMs trained on public 

\section{Problem statement}\label{sec:problem}

The background discussion demonstrates the clear problems within this domain, it has proven that cyber security labs and AI tutors give real value to the end user. But the problem lies in whether they can be combined into a coherent, safe, and educationally useful system that generates individuals vulnerable, isolated environments on demand, then expose them through controlled infrastructure. From here, it should support the learner with contextual guidance, and reliably clean-up the environment afterwards. Put simply, VulnForge will be investigating whether generative AI can move from just being an accessory that people use, to being part of the content production and delivery pipeline itself. 

This leads to the following questions:
\begin{enumerate}
    \item To what extent can a \gls{llm} pipeline generate vulnerability specifications that are both deployable and sufficiently diverse to support cyber security training within defined challenge categories?
    \item How can generated challenges be deployed and exposed through isolated, short-lived environments without undermining host safety , challenge integrity or resource control?
    \item Relative to a purely curated challenge library, does an on-demand generated approach offer meaningful improvement in re-playability and education?
\end{enumerate}